% 12 的正因数在整除关系下的偏序（Hasse 图）。节点旁为 Möbius 函数 mu(1,x)：
% mu(1,1)=1；mu(1,2)=mu(1,3)=-1；mu(1,4)=-(1-1)=0；mu(1,6)=-(1-1-1)=1；mu(1,12)=-(1-1-1+0+1)=0。
\begin{tikzpicture}[line width=.8pt,font=\figurefont\fontsize{8}{10}\selectfont,
  vtx/.style={circle,draw=FigureStroke,fill=FigureBlueFill,minimum size=7mm,inner sep=0pt,line width=.8pt},
  mu/.style={FigureInk,font=\figurefont\fontsize{8}{10}\selectfont}]
  \node[vtx,fill=FigureYellowFill] (1) at (0,0) {$1$};
  \node[vtx] (2) at (-1.1,1.25) {$2$};
  \node[vtx] (3) at (1.1,1.25) {$3$};
  \node[vtx] (4) at (-2.2,2.5) {$4$};
  \node[vtx] (6) at (0,2.5) {$6$};
  \node[vtx] (12) at (-1.1,3.75) {$12$};
  \foreach \a/\b in {1/2,1/3,2/4,2/6,3/6,4/12,6/12}{\draw[FigureStroke] (\a)--(\b);}
  \node[mu,right=2pt of 1] {$\mu(1,1)=1$};
  \node[mu,left=2pt of 2] {$-1$};
  \node[mu,right=2pt of 3] {$-1$};
  \node[mu,left=2pt of 4] {$0$};
  \node[mu,right=2pt of 6] {$1$};
  \node[mu,right=2pt of 12] {$0$};
\end{tikzpicture}
